% from ega0-22-fr.tex
\oldpage[0\textsubscript{IV-1}]{191}

\subsection{Application aux relations entre certains anneaux locaux et leurs complétés}
\label{subsection:0.22.3-fr}

\begin{lemma}[22.3.1]
\label{0.22.3.1-fr}
Soient $A_0$ un anneau semi-local noethérien, $A$ une $A_0$-algèbre finie
(qui est donc un anneau semi-local noethérien, cf. Bourbaki,
\emph{Alg. comm.}, chap.~IV, §~2, n\textsuperscript{o}~5, cor.~3 de la
prop.~9). Si $\widehat A_0$ et $\widehat A$ sont les complétés de $A_0$ et
$A$ pour leurs topologies préadiques respectives, alors, lorsqu'on munit
$A_0$, $A$ et $\widehat A$ des topologies discrètes, $\widehat A$ est une
$A$-algèbre formellement lisse relativement à $A_0$ (19.9.1).
\end{lemma}

\begin{proof}
On sait en effet que
$\widehat A=A\otimes_{A_0}\widehat A_0$ (Bourbaki,
\emph{Alg. comm.}, chap.~IV, §~2, n\textsuperscript{o}~5, cor.~3 de la
prop.~9 et chap.~III, §~3, n\textsuperscript{o}~4, th.~1), donc la proposition
résulte de (19.9.3).
\end{proof}

\begin{proposition}[22.3.2]
\label{0.22.3.2-fr}
Soient $A$ un anneau local noethérien régulier, $K$ son corps des fractions,
$p$ la caractéristique de $K$. On suppose vérifiée l'une des hypothèses
suivantes~:
\begin{enumerate}
  \item[(i)] $p=0$.
  \item[(ii)] $p>0$ et il existe une famille filtrante décroissante
  $(A_\alpha)$ de sous-anneaux noethériens de $A$, tels que $A$ soit une
  $A_\alpha$-algèbre finie pour tout $\alpha$, que
  $A^{p^{h(\alpha)}}\subset A_\alpha$ pour un entier $h(\alpha)$ convenable,
  et que, si $K_\alpha$ désigne le corps des fractions de $A_\alpha$, on ait
  \[
    \bigcap_\alpha K_\alpha(K^p)=K^p.
  \]
\end{enumerate}
Soient alors $B$ une $A$-algèbre finie intègre et contenant $A$,
$\mathfrak n$ un idéal premier de $B$, $C$ l'anneau local $B_\mathfrak n$,
$\mathfrak q$ un idéal premier du complété $\widehat C$ tel que
$\mathfrak q\cap C=0$, de sorte que l'anneau local
$\widehat C_\mathfrak q$ est une algèbre sur le corps des fractions $L$ de
$B$. Alors $\widehat C_\mathfrak q$ est une $L$-algèbre formellement lisse
pour sa topologie $\mathfrak q$-adique, et par suite un anneau
géométriquement régulier sur $L$ (19.6.5).
\end{proposition}

\begin{proof}
Distinguons deux cas.

I) Supposons que $\mathfrak n$ contienne l'idéal maximal $\mathfrak m$ de
$A$, et par suite $\mathfrak n\cap A=\mathfrak m$. Alors $\mathfrak n$ est
maximal (Bourbaki, \emph{Alg. comm.}, chap.~V, §~2,
n\textsuperscript{o}~1, prop.~1) et si $\mathfrak r$ est le radical de
l'anneau semi-local $B$, $\widehat C$ est le séparé complété de $B$ pour la
topologie $\mathfrak n$-préadique, et un des composants de l'anneau
semi-local $\widehat B$, complété de $B$ pour la topologie
$\mathfrak r$-préadique (Bourbaki, \emph{Alg. comm.}, chap.~III, §~3,
n\textsuperscript{o}~4, prop.~8)~; on peut donc écrire
$\widehat C_\mathfrak q=\widehat B_{\mathfrak q'}$, où $\mathfrak q'$ est
l'idéal premier de $\widehat B$ image réciproque de $\mathfrak q$. Notons
maintenant que $B$ est un sous-anneau de $\widehat B$ et l'hypothèse
$\mathfrak q\cap C=0$ entraîne $\mathfrak q'\cap B=0$, donc
$\widehat B_{\mathfrak q'}$ est aussi localisé de l'anneau
$\widehat B\otimes_B L$ en un de ses idéaux premiers $\mathfrak q''$, dont
$\mathfrak q'$ est l'image réciproque dans $\widehat B$. En outre on a
$\widehat B=\widehat A\otimes_A B$ (Bourbaki, \emph{Alg. comm.}, chap.~IV,
§~2, n\textsuperscript{o}~5, cor. de la prop.~9 et chap.~III, §~3,
n\textsuperscript{o}~4, th.~1), donc
$\widehat B\otimes_B L=\widehat A\otimes_A L$~; enfin, si $\mathfrak p$ est
l'idéal premier de $\widehat A$, image réciproque de $\mathfrak q''$, le
fait que $A\to B$ est injectif et la commutativité du diagramme
\[
  \xymatrix{
    B \ar[r] & \widehat B\\
    A \ar[u] \ar[r] & \widehat A \ar[u]
  }
\]
entraînent que $\mathfrak p\cap A=0$, donc
$(\widehat A\otimes_A L)_{\mathfrak q''}$ est aussi localisé de l'anneau
\[
  \widehat A_\mathfrak p\otimes_A L
  =(\widehat A_\mathfrak p\otimes_A K)\otimes_K L
\]
en un de ses idéaux premiers.

